\chapter{Introdução}
\label{cap:introducao}

Compreender como o risco é remunerado nos mercados financeiros é um dos objetivos
centrais da teoria moderna de finanças. Quando os retornos de centenas de ativos se
movem de forma coordenada em resposta a choques macroeconômicos, variações nas taxas
de juros ou mudanças no apetite global por risco, essa covariância reflete a exposição
comum dos ativos a fontes de risco que atravessam o mercado como um todo. Separar esse
componente sistemático da variação genuinamente específica de cada empresa, o componente
idiossincrático, é o problema que os modelos de fatores se propõem a resolver. A
relevância prática é direta: estimação eficiente de matrizes de covariância, construção
de portfólios diversificados e formação de expectativas condicionais de retorno dependem
de uma descrição adequada dessa estrutura fatorial subjacente.

A primeira resposta a esse problema veio com o CAPM de \textcite{sharpe1964}, que
forneceu um arcabouço teórico simples e testável: o retorno esperado de cada ativo é
determinado por sua sensibilidade a um único fator sistemático, o excesso de retorno do
portfólio de mercado. \textcite{ross1976} generalizou essa estrutura com a
\textit{Arbitrage Pricing Theory}, admitindo múltiplos fatores sem exigir hipóteses de
equilíbrio. Sob ausência de arbitragem, os retornos esperados são determinados por
sensibilidades a um vetor de fatores sistemáticos, cada um carregando seu próprio prêmio
de risco. Em ambos os casos, as exposições são estimadas por regressão de série temporal
e tratadas como constantes ao longo do tempo.

À medida que novas evidências empíricas surgiram, as limitações do modelo de fator único
tornaram-se evidentes. \textcite{jensen1968} documentou desempenhos anômalos não
explicados pelo CAPM, e \textcite{banz1981} identificou o ``efeito tamanho'', segundo o
qual ações de menor capitalização apresentam retornos sistematicamente superiores ao
previsto pelo modelo. Esses resultados indicaram a existência de outras dimensões de
risco sistemático não capturadas por um único fator. A resposta foi a segunda geração de
modelos, que mantém a estrutura linear mas constrói os fatores diretamente a partir de
características observáveis das empresas: \textcite{fama1993} propõem portfólios formados
com base em tamanho e índice book-to-market como proxies para fontes de risco não
capturadas pelo mercado; \textcite{fama2015} estendem o modelo para cinco fatores,
incorporando rentabilidade e investimento; \textcite{carhart1997} acrescenta um fator de
momento. A lógica permanece a mesma: hipóteses ex-ante sobre quais características geram
prêmios sistemáticos, estimação por regressão linear de painel e exposições constantes no
tempo.

Três limitações estruturais acumuladas ao longo dessa trajetória motivaram a busca por
alternativas. A primeira é a dependência da especificação ex-ante. O modelo é tão bom
quanto as hipóteses do pesquisador sobre quais características são portadoras de prêmio
de risco, e a proliferação de centenas de fatores documentados sem critério de seleção
sistemático, o chamado \textit{factor zoo} \parencite{harvey2016}, expôs o caráter
arbitrário dessa escolha. A segunda é a restrição linear, que impede a captura de
interações não-lineares entre características e retornos. A terceira é a estacionariedade
das exposições: betas constantes no tempo desconsideram que a sensibilidade de uma
empresa ao câmbio, ao ciclo de commodities ou às condições de crédito se altera com o
ambiente macroeconômico.

A literatura recente passou a concentrar esforços em modelos que deixam os dados
determinar os fatores, dispensando a especificação ex-ante. \textcite{kelly2019} propõem
o IPCA como caminho que preserva a tratabilidade econométrica linear e ao mesmo tempo
admite fatores latentes com exposições variáveis no tempo: as exposições passam a ser
função linear das características dos ativos, e os fatores são extraídos por uma versão
instrumentada de PCA. No teste empírico no mercado norte-americano, cinco fatores IPCA
explicam a seção transversal de retornos médios com precisão substancialmente superior à
de modelos com fatores observáveis, e apenas dez de trinta e seis características
testadas se mantêm estatisticamente significativas a 1\%, indicando ganho de parcimônia
além do ganho de ajuste. O resultado é relevante para este trabalho porque demonstra que
o trade-off entre flexibilidade e disciplina econométrica pode ser resolvido sem
abandonar a estrutura linear, desde que se admitam fatores latentes condicionados em
características.

\textcite{gu2021} avançam ao remover também a restrição de linearidade. Os autores
propõem o \textit{Conditional Autoencoder}, no qual fatores latentes e exposições
condicionadas em características são aprendidos conjuntamente por uma rede neural,
permitindo que interações não-lineares entre características e retornos sejam capturadas
sem especificação ex-ante. A comparação direta com IPCA e com modelos lineares de
Fama-French no \textit{cross-section} norte-americano mostra ganhos sistemáticos de
$R^2$ fora da amostra e melhor performance de portfólios construídos a partir das
previsões do modelo. O trabalho fornece a base arquitetural sobre a qual modelos neurais
subsequentes são construídos: extrator de características, fatores latentes condicionados
e cálculo dos retornos como combinação linear dos fatores ponderada por exposições.

\textcite{duan2022} identificam uma limitação que persiste mesmo nos modelos neurais
determinísticos. Os retornos de ações apresentam razão sinal-ruído extremamente baixa, o
que dificulta o aprendizado dos fatores latentes ainda que a arquitetura admita
não-linearidades. A proposta dos autores, o \textit{FactorVAE}, ataca esse problema com
duas mudanças principais. A primeira é tratar os fatores latentes como variáveis
aleatórias e não como pontos no espaço latente, recuperando o tratamento probabilístico
das fontes de risco característico dos modelos clássicos de fatores e regularizando o
aprendizado em dados ruidosos. A segunda, e mais importante na prática, é um esquema de
treinamento professor-aluno: durante o treinamento, um módulo com acesso aos retornos
contemporâneos extrai os fatores que melhor reconstroem aqueles retornos e funciona como
alvo para um segundo módulo, treinado exclusivamente com informação histórica disponível
na data da previsão. O modelo aprende não apenas a prever retornos, mas a prever os
fatores que um observador informado teria identificado, o que se traduz em ganho
empírico documentado pelos autores no mercado A-share chinês: superioridade consistente
sobre IPCA, sobre o \textit{Conditional Autoencoder} e sobre modelos preditivos baseados
em redes recorrentes, tanto em \textit{Information Coefficient} cross-sectional quanto em
performance de portfólio.

A avaliação original do \textit{FactorVAE} foi conduzida em um universo superior a 3.000
ativos com período de teste de dois anos. Esse recorte deixa em aberto se a vantagem do
modelo sobrevive em mercados estruturalmente menores e em janelas de teste mais longas.
A pergunta que organiza este trabalho é se o \textit{FactorVAE}, replicado fielmente no
mercado brasileiro de ações, mantém vantagem de \textit{skill} cross-sectional sobre
benchmarks lineares e neurais simples, e se essa vantagem sobrevive à tradução em
portfólio sob custos de transação realistas. O contexto brasileiro impõe condições
substancialmente diferentes do estudo original. O universo de ativos negociados na B3 é
uma ordem de grandeza menor que o A-share chinês, apresenta marcada concentração
setorial em commodities, energia e setor financeiro, e exibe liquidez heterogênea entre
os tickers. O período de teste, por sua vez, cobre sete anos, contra dois anos do
estudo original.

Este trabalho faz três contribuições. A primeira é documentar a replicação do
\textit{FactorVAE} fora do mercado chinês, registrando as adaptações necessárias ao
contexto da B3: construção do universo \textit{point-in-time} com filtro mínimo de
liquidez, vinte características técnicas derivadas de preço e volume e hiperparâmetros
reportados integralmente. A segunda é conduzir uma avaliação sistemática em duas
dimensões: \textit{skill} cross-sectional, medido por Rank IC e Rank ICIR, e performance
de portfólio após custos de transação, medida por retorno anualizado, Sharpe e drawdown
máximo, comparando o \textit{FactorVAE} com três modelos de referência em progressão
deliberada de complexidade: GRU, IPCA e \textit{Conditional Autoencoder}. A terceira é
confrontar os resultados diretamente com os de \textcite{duan2022}, identificando as
dimensões em que a replicação confirma o artigo original, as que divergem em magnitude,
e as hipóteses para as discordâncias observadas.

O restante do trabalho está organizado da seguinte forma. O Capítulo~\ref{cap:revisao}
percorre a evolução dos modelos de fatores, partindo dos modelos estáticos e chegando aos
modelos com fatores latentes dinâmicos, estabelece os fundamentos de inferência
variacional necessários para compreender a função objetivo do modelo e apresenta o
\textit{FactorVAE} em detalhe conceitual. O Capítulo~\ref{cap:metodologia} especifica
operacionalmente a arquitetura, os dados utilizados, os modelos de comparação e os
procedimentos de avaliação. O Capítulo~\ref{cap:resultados} apresenta e interpreta os
resultados nas quatro dimensões de avaliação. O Capítulo~\ref{cap:conclusao} sintetiza os
achados, responde à pergunta de pesquisa, registra as limitações e propõe extensões.